%% Section 8 — Conclusion
\section{Conclusion}
\label{sec:conclusion}

This paper began from a bibliometric fact: nearly five thousand publications on data centers, and not one engaging the theoretical frameworks through which innovation scholars analyse infrastructure transitions. The argument it has developed is that this gap is not incidental but produced — by a metaphor that naturalises the infrastructure, by a disciplinary division of labour that assigns it to engineers and ethnographers, and by the deliberate opacity that operators have constructed around it. The cloud conceals.

The concept of \textit{datacentering} is proposed as a remedy. The shift from noun to gerund is a shift in what we attend to: from the stable properties of a technology to the active, contested, ongoing process of assembling it. When we study datacentering rather than data centers, we immediately encounter the features that innovation scholars are trained to analyse — actors, governance, contestation, imaginaries, regime dynamics — and we immediately face the methodological challenge that those features present: they are constituted across multiple registers (textual, visual, financial, temporal), distributed across multiple scales (facility, city, jurisdiction, global network), and partially hidden by design.

The Complex Adaptive Systems framing, drawn from \citet{phillips2019complex} as a sensemaking lens rather than a formal model, organises the methodological response. The conceptual dimension of the assemblage — how actors construct imaginaries and manage visibility — is addressed through corporate and territorial text analysis (WP1) and visual and satellite imagery (WP2). The structural dimension — the ownership hierarchy and actor constellation through which governance capacity is constituted or dissolved — is addressed through ownership cartography (WP3) and urban-embeddedness analysis. The temporal dimension — how contestation unfolds and what determines whether regime actors respond — is addressed through contestation signal detection (WP4). Together, the four work packages constitute a multimodal cartographic methodology that is adequate to the CAS-like complexity of the object, and the Virginia--Amsterdam comparison demonstrates, illustratively, the analytical purchase of applying it.

The conclusion the paper reaches is that no single researcher or team can build this cartography alone. The geography is too large, the data too distributed, the opacity too deliberate, and the jurisdictional knowledge too local for any single team to map it adequately. This is not a counsel of humility but a methodological finding: the research commons is the only architecture adequate to the object's actual complexity. The open facility index, the shared analytical pipelines, the contributing protocols, and the GitHub Discussions are not supplementary to the scholarly contribution — they are part of it.

We therefore close with an explicit invitation. If you have local knowledge about a planning dispute, a facility that does not appear in published records, a moratorium under negotiation, or a community campaign in its early stages — that knowledge is analytically valuable and there is a structured place for it in this research infrastructure. If you have methodological expertise in NLP, computer vision, network analysis, or event detection — the pipelines are open and documented. If you have access to national company registries, REIT filings, or energy and water disclosure records that are not publicly accessible through the channels the project currently uses — get in touch.

The cloud is not a metaphor. It is infrastructure. And infrastructure, with the right tools and the right collaborators, can be mapped.
